\part{Physical design}\label{tm:part:physical-design}

\chapter{Flow overview}\label{tm:ch:flow-overview}

This chapter describes the physical design flow. It shows the steps from netlist to layout.

\chapter{Constraint files}\label{tm:ch:constraint-files}

This chapter describes the constraint files of the flows.

\chapter{Sky130 flow}\label{tm:ch:sky130-flow}

This chapter describes the Sky130 flow. It gives the configuration and the results.

\chapter{ASAP7 flow}\label{tm:ch:asap7-flow}

This chapter describes the ASAP7 flow. It gives the configuration and the results.

\chapter{FreePDK45 flow}\label{tm:ch:freepdk45-flow}

This chapter describes the FreePDK45 flow. It gives the configuration and the status.

\chapter{Commercial flow comparison}\label{tm:ch:commercial-flow-comparison}

This chapter compares the open flow with the commercial flows for one macro.

% Moved verbatim from CLAUDE.md by tools/docs/md2tex.py; edit this file, not the original.
\chapter{Technology Node Strategy}\label{tm:pd:node-strategy}
\begin{relocated}{CLAUDE.md}

Three nodes are supported by the existing \src{pnr/Makefile}. Use the
right node for the right goal:

{\footnotesize\setlength{\tabcolsep}{3pt}

{
\begin{longtable}[]{@{}
  >{\raggedright\arraybackslash}p{\dimexpr(\linewidth-10\tabcolsep)*111/1000\relax}
  >{\raggedright\arraybackslash}p{\dimexpr(\linewidth-10\tabcolsep)*214/1000\relax}
  >{\raggedright\arraybackslash}p{\dimexpr(\linewidth-10\tabcolsep)*214/1000\relax}
  >{\raggedright\arraybackslash}p{\dimexpr(\linewidth-10\tabcolsep)*111/1000\relax}
  >{\raggedright\arraybackslash}p{\dimexpr(\linewidth-10\tabcolsep)*83/1000\relax}
  >{\raggedright\arraybackslash}p{\dimexpr(\linewidth-10\tabcolsep)*246/1000\relax}@{}}
\toprule\noalign{}
\begin{minipage}[b]{\linewidth}\raggedright
Node
\end{minipage} & \begin{minipage}[b]{\linewidth}\raggedright
Makefile target
\end{minipage} & \begin{minipage}[b]{\linewidth}\raggedright
Config status
\end{minipage} & \begin{minipage}[b]{\linewidth}\raggedright
Realistic CPU freq
\end{minipage} & \begin{minipage}[b]{\linewidth}\raggedright
Fabricatable
\end{minipage} & \begin{minipage}[b]{\linewidth}\raggedright
Best use
\end{minipage} \\
\midrule\noalign{}
\endhead
\bottomrule\noalign{}
\endlastfoot
\textbf{Sky130 130nm} & \code{librelane-sky130-cpu},
\code{librelane-sky130-soc} & \ok{} Complete (\src{pnr/sky130/}) &
75--100 MHz & \ok{} Yes (SkyWater MPW) & Functional sign-off; tape-out
candidate \\
\textbf{FreePDK45 / NanGate45 45nm} & \code{librelane-nangate45} & \ok{}
Complete (\src{pnr/freepdk45/}) & 150--250 MHz & No (predictive) & First
non-Sky130 exploration; SRAM macros available \\
\textbf{ASAP7 7nm FinFET} & \code{librelane-asap7} & \ok{} Complete
(\src{pnr/asap7/}); Run 43 signoff 2026-05-20 & \textbf{1418 MHz
achieved (Run 43)} & No (predictive) & Phase 2+3 RTL PPA sign-off
completed at 1418 MHz / 27.27 mW / 3 844 µm² \\
\end{longtable}
}

}

\section{\texorpdfstring{Commercial flows (\src{pnr/commercial/}, Sky130
only, indicative)}{Commercial flows (, Sky130 only, indicative)}}\label{tm:pd:node-strategy:commercial-flows-pnrcommercial-sky130-only-indicative}

DC -\textgreater{} ICC2 and Genus -\textgreater{} Innovus flows for the
Sky130 CPU macro \code{rv32i_cpu_top}, built on the EDA host (no nix).
They are judged with OpenROAD B2 by one StarRC -\textgreater{}
PrimeTime/\allowbreak{}PrimePower chain, plus
KLayout/\allowbreak{}Magic/\allowbreak{}Netgen,
Formality/\allowbreak{}Conformal and RedHawk. Final arms are ICC2 b6 and
Innovus b6. All figures are indicative: StarRC is calibrated to OpenRCX,
not silicon, and the SRAM Liberty is TT-only (\code{o1i}). No arm meets
sign-off. Results and caveats are in
\src{docs/COMMERCIAL_FLOW_COMPARISON.md}.

\section{Legacy P\&R directories (superseded --- do not start new work
here)}\label{tm:pd:node-strategy:legacy-pr-directories-superseded-do-not-start-new-work-here}

\src{pnr/librelane/} and \src{pnr/openlane1/} are the pre-per-node
Sky130 CPU flows. Both are \textbf{superseded by \src{pnr/sky130/cpu/}
and \src{pnr/sky130/soc/}}, which are the live Sky130 flows and the only
ones with real Magic DRC + Netgen LVS enabled.
\src{pnr/sky130/cpu/config.json} was forked from
\src{pnr/librelane/config.json} (see
\src{docs/SKY130_REAL_DRC_LVS_EVALUATION.md} §Stage 1). Their targets
(\code{librelane-sky130},
\code{docker-run}/\code{docker-synth}/\code{docker-results}) still exist
as a fallback, but neither directory has run since 2026-03-25 and both
carry the stale \code{VERILOG_FILES} package-closure gap tracked in bead
\bead{135} --- run
\texttt{python3\ tools/\allowbreak{}verif/\allowbreak{}check\_\allowbreak{}source\_\allowbreak{}closure.\allowbreak{}py\ <config>}
before trusting either. \code{pnr/openlane/} was speculative OpenLane-2
scaffolding that was never wired to any Makefile target and has been
deleted (LibreLane is the OL2-lineage tool this project actually uses).
\src{pnr/freepdk45/} is \textbf{not} superseded --- it is a
dormant-but-intentional third node with no per-node replacement, and
\code{librelane-nangate45} remains the only path to it.

\section{Sky130 real DRC/LVS sign-off (GH epic \#102, Stages 1--2
complete 2026-07-31)}\label{tm:pd:node-strategy:sky130-real-drclvs-sign-off-gh-epic-102-stages-12-complete-2026-07-31}

Sky130 is the \textbf{only} node in this project with genuine physical
verification --- ASAP7 skips
Magic/\allowbreak{}KLayout/\allowbreak{}Netgen entirely (predictive PDK)
and FreePDK45 has them permanently blocked.

\begin{itemize}
\tightlist
\item
  \textbf{Stage 1} --- CPU standalone (\src{pnr/sky130/cpu/}): KLayout
  DRC 0, Netgen LVS MATCH, +0.366 ns setup @ 75 MHz. Magic DRC 27.7 M is
  a waived PDK artifact --- 100 \% of violations sit inside SRAM macro
  footprints (bead \bead{45a}).
\item
  \textbf{Stage 2} --- CPU macro + peripherals SoC, no GPU
  (\src{pnr/sky130/soc/}): LVS PASSED, routing DRC 0, PDN 0, hold clean,
  40 MHz, nom\_tt 37.63 mW, die 6700 × 3100 µm. Residual gates, neither
  open work: ss setup fail (\code{ujv}, deferred host-blocked), antenna
  140/120 (\code{58q}, closed as an accepted trade-off against hold).
\item
  \warn{} \textbf{Stage-2 timing is typical-corner (nom\_tt) only.} The
  SRAM macro is characterized at TT alone, so "clean at 9 corners" means
  9 corner \emph{labels} backed by one macro model. Physical gates
  (LVS/\allowbreak{}DRC/\allowbreak{}PDN) are unaffected. Per-corner
  SRAM characterization is host-blocked (\textasciitilde80--95 h SPICE;
  bead \bead{o1i}).
\item
  \textbf{Stages 3--4} (GPU, \ghref{105}/\allowbreak{}\#106): host-gated on
  ≥32 GB RAM + \textasciitilde500 GB scratch. GPU stays
  ASAP7-indicative.
\item
  Detail → \src{docs/SKY130_REAL_DRC_LVS_EVALUATION.md} §7.
\end{itemize}

\section{Sky130 PPA ceiling (assessed 2026-03-28)}\label{tm:pd:node-strategy:sky130-ppa-ceiling-assessed-2026-03-28}

Sky130 130nm is frequency-limited by cell drive strength and clock tree
skew. The realistic ceiling for this pipeline topology is
\textbf{\textasciitilde100--120 MHz} with all optimisations applied ---
not the 200--500 MHz originally targeted.

{
\begin{longtable}[]{@{}
  >{\raggedright\arraybackslash}p{\dimexpr(\linewidth-4\tabcolsep)*517/1000\relax}
  >{\raggedright\arraybackslash}p{\dimexpr(\linewidth-4\tabcolsep)*127/1000\relax}
  >{\raggedright\arraybackslash}p{\dimexpr(\linewidth-4\tabcolsep)*336/1000\relax}@{}}
\toprule\noalign{}
\begin{minipage}[b]{\linewidth}\raggedright
Technique
\end{minipage} & \begin{minipage}[b]{\linewidth}\raggedright
Gain
\end{minipage} & \begin{minipage}[b]{\linewidth}\raggedright
Effort
\end{minipage} \\
\midrule\noalign{}
\endhead
\bottomrule\noalign{}
\endlastfoot
Switch to \code{sky130_fd_sc_hs} (high-speed cell library) & +10--30 MHz
& Low --- Makefile config change \\
EX-stage retiming (split ALU + branch-compare) & +15--25 MHz & Medium
--- re-verify forwarding paths \\
SDC multi-cycle exemptions for CSR/debug paths & +5--10 MHz & Low ---
SDC annotation only \\
Floorplan: place forwarding mux adjacent to ALU & +5--10 MHz & Medium
--- P\&R iteration \\
\textbf{Combined ceiling} & \textbf{\textasciitilde100--120 MHz} & \\
\end{longtable}
}

\section{\texorpdfstring{FreePDK45/\allowbreak{}NanGate45 (ready to
run)}{FreePDK45/NanGate45 (ready to run)}}\label{tm:pd:node-strategy:freepdk45nangate45-ready-to-run}

\src{pnr/freepdk45/} is fully configured: \code{config.json} (2.5 ns /
400 MHz target), SRAM macros (\code{sram_1rw_256x32_freepdk45} with
LEF/\allowbreak{}LIB/\allowbreak{}GDS), \code{macro_placement.cfg},
\code{pdn.tcl}, \code{freepdk45.sdc}, \code{freepdk45.upf}. Run with
\texttt{make\ librelane-\allowbreak{}nangate45}. Expected Phase 2
result: 150--250 MHz, \textasciitilde8--10× logic density vs Sky130.

\section{ASAP7 (Phase 2+3 sign-off achieved at Run 43, 2026-05-20)}\label{tm:pd:node-strategy:asap7-phase-23-sign-off-achieved-at-run-43-2026-05-20}

\src{pnr/asap7/} is fully configured and signed off: \code{config.json}
(705 ps / 1418 MHz, CTS clustering 8/10, density 50 \%),
\code{constraints/asap7.sdc} (single-period clock,
hold/\allowbreak{}setup uncertainty 10/15 ps, SRAM ICG \code{gclk_sram}
generated clock + false-paths), \code{macro_placement.cfg},
\code{pdn.tcl}, \code{sram_1rw_256x32_asap7_TT_0p7V_25C.lib} (RVT TT 0.7
V / 25 °C). Run with \texttt{make\ librelane-\allowbreak{}asap7}. Phase
2+3 PPA (no perf counters): \textbf{1418 MHz fmax, 27.27 mW power, 3 844
µm² stdcell, 0 DRC, 0 antenna \warn{} unsupported --- see bead
\bead{xy6}, +5.97 ps setup slack, +22.54 ps hold slack} at run
\code{pnr/asap7/runs/RUN_2026-05-20_06-27-10/} (run directory no longer
on disk to re-verify directly; \code{drt.tcl}\textquotesingle s
catch-and-continue patch was already in place before this run per
\src{docs/ASAP7_RUN_HISTORY.md}\textquotesingle s caveat banner, so
timing/\allowbreak{}power/\allowbreak{}area stand but the
DRC/\allowbreak{}antenna claim does not). \textbf{Phase 5 M7 update
(2026-06-02):} adding the perf-monitoring CSRs lowers the standalone CPU
sign-off to \textbf{1282 MHz (780 ps, +24 ps setup, +28 ps hold, 24.38
mW, 4 058 µm²; achievable \textasciitilde1323 MHz)} at run
\code{pnr/asap7/cpu/runs/RUN_2026-06-02_20-37-12/} --- the 1418 MHz
number does not hold with M7 (limiter = 836-fanout counter-write-enable
cone). Moot for the SoC (GPU-governed \textasciitilde571 MHz).
\src{pnr/asap7/cpu/config.json} + \code{constraints/asap7.sdc} are now
at 780 ps. See \src{docs/ASAP7_RUN_HISTORY.md} for the full campaign
history. \warn{} \textbf{Both the 1418 MHz and 1282 MHz figures above
are \code{8f3}-affected} (bead \bead{8f3}/\code{0p6}, 2026-09-16): stock
\code{OpenROAD.STAMidPNR-3} without an in-session re-route, same
zero-wire STA defect as the SoC. Run 43\textquotesingle s own artifacts
are gone (unrecoverable); the M7 1282 MHz point \textbf{was re-measured
on current RTL with the fix} (\code{STA_POSTGRT_INSESSION_GRT=1}, run
tag \code{bpp7_postgrt_set}) and does \textbf{not} hold --- honest
result is setup WNS \textbf{−169.7 ps / TNS −70,913 ps / 1008 violators}
(hold clean, unaffected). Full detail: \src{docs/ASAP7_RUN_HISTORY.md}.
\end{relocated}
